\documentclass[a4paper]{article}
\usepackage{fullpage}
\usepackage[margin=63pt]{geometry}
\usepackage{graphics}
\usepackage{xcolor}
\usepackage{graphicx}
\usepackage{subcaption}
\usepackage{amssymb}
\usepackage{pdfpages}
\usepackage{booktabs}

\input{../stachnisslab-math.tex}
%% Standard latex definitions used at the Stachniss-Lab

\usepackage{graphics}           
\usepackage{times}              
\usepackage{amsmath}            
\usepackage{amssymb}            
\usepackage{graphicx}
\usepackage{algorithm}
\usepackage[noend]{algpseudocode}
\usepackage{booktabs}
\usepackage{balance}
\usepackage{color}
\usepackage[nocompress]{cite} % sorts citations automatically, e.g., "[1],[2],[3]" instead of "[2],[3],[1]".
\definecolor{instructioncolor}{rgb}{.5,.5,.5}
\usepackage{multirow} 
\usepackage{url} 
\newcommand{\instruction}[1]{\textcolor{instructioncolor}{#1}}

% make caption font small for better separation of figures and text
\usepackage[font=small]{caption}
\usepackage{svg}
 %% Key definitions for text elements. USE ONLY THEM! Do not use naked \ref{}.
\def\secref#1{Sec.~\ref{#1}}
\def\figref#1{Fig.~\ref{#1}}
\def\tabref#1{Tab.~\ref{#1}}
\def\eqref#1{Eq.~(\ref{#1})}
\def\algref#1{Alg.~\ref{#1}}

%% Other useful macros
\newcommand\todo[1]{\textbf{[TODO: #1}]}

\makeatletter
\usepackage{xspace}
%% the \onedot macro is producing only one dot at line ends.
%% thus \etal. will not produce et al..
\DeclareRobustCommand\onedot{\futurelet\@let@token\@onedot}
\def\@onedot{\ifx\@let@token.\else.\null\fi\xspace}
\def\eg{e.g\onedot} \def\Eg{E.g\onedot}
\def\ie{i.e\onedot} \def\Ie{I.e\onedot}
\def\cf{cf\onedot} \def\Cf{Cf\onedot}
\def\etc{etc\onedot} \def\vs{vs\onedot}
\def\wrt{w.r.t\onedot} \def\dof{d.o.f\onedot}
%% Cyrill does not like emph...
%\def\etal{\emph{et al}\onedot}
\def\etal{{et al}\onedot}
\makeatother

\def\etalcite#1{\etal~\cite{#1}}

\usepackage{array}
%% this allows to use something like p{2cm} as column type, but with left, center, and right alignment
\newcolumntype{L}[1]{>{\raggedright\let\newline\\\arraybackslash\hspace{0pt}}m{#1}}
\newcolumntype{C}[1]{>{\centering\let\newline\\\arraybackslash\hspace{0pt}}m{#1}}
\newcolumntype{R}[1]{>{\raggedleft\let\newline\\\arraybackslash\hspace{0pt}}m{#1}}




\def\etalcite#1{\etal~\cite{#1}}
\newcommand\circledmark[1][red]{%
  \ooalign{%
    \hidewidth
    \kern0.65ex\raisebox{-0.9ex}{\scalebox{3}{\textcolor{#1}{\textbullet}}}
    \hidewidth\cr
    $\checkmark$\cr
  }%
}

\definecolor{brewerGreen0}{HTML}{E5F5F9}
\definecolor{brewerGreen1}{HTML}{99D8C9}
\definecolor{brewerGreen2}{HTML}{2CA25F}
\definecolor{brewerCyan0}{HTML}{ECE2F0}
\definecolor{brewerCyan1}{HTML}{A6BDDB}
\definecolor{brewerCyan2}{HTML}{1C9099}
\definecolor{brewerCyan9}{HTML}{023858}
\definecolor{brewerGrey0}{HTML}{F0F0F0}
\definecolor{brewerGrey1}{HTML}{BDBDBD}
\definecolor{brewerGrey2}{HTML}{636363}

\definecolor{commentColor}{HTML}{000000}
\definecolor{revisionColor}{HTML}{0238A8}
\definecolor{answerBoxColor}{HTML}{006D2C}

\usepackage[colorlinks=true,linkcolor=answerBoxColor,citecolor=answerBoxColor,urlcolor=answerBoxColor,pagebackref]{hyperref}

\usepackage{tcolorbox}
\def\answer#1#2{
  \begin{tcolorbox}[width=\textwidth,colback={white},title={\textbf{Answer
    #1}},colbacktitle=answerBoxColor,coltitle=white] #2
\end{tcolorbox} }

\newcounter{todonum}
\newcommand\resettodonum{\setcounter{todonum}{0}}
\newcommand\mytodo{%
  \stepcounter{todonum}%
  \textcolor{red}{\;TODO \thetodonum}
}

\newcommand\boldblue[1]{\textcolor{blue}{\textbf{#1}}}
\newcommand\boldred[1]{\textcolor{red}{\textbf{#1}}}
\newcommand{\reviewed}[1]{\textcolor{blue}{#1}}

\title{
  \textbf{Semantic-Landmark Particle Filter for Reliable Robot Localisation in Vineyards}
  \\[\baselineskip] Authors' Response to ICRA 2026 AE and Reviewers
}
\author{Anonymous Authors}

\usepackage{titling}
\setlength{\droptitle}{-0.5in}

\begin{document}
\maketitle
\resettodonum

We thank the Associate Editor and both reviewers for their careful and constructive feedback.
In the revised manuscript, we addressed all requested points and added new experiments for fairness,
robustness, and runtime reporting.

\paragraph{Summary of major revisions}
\begin{itemize}
  \item \textbf{Fairer benchmarking:} added \textbf{AMCL+NoisyGPS} (Kalman fusion) baseline in both traverses and reported multiseed results (seeds 11, 22, 33).
  \item \textbf{Clearer metrics:} report raw APE and \textbf{RPE at 2 m / 5 m} explicitly, plus cross-track error, row-correct fraction, and row-switch events.
  \item \textbf{Semantic-wall clarity:} clarified the methodology and equations for ray casting and semantic-wall likelihood terms, and added a point-only non-wall formulation for comparison.
  \item \textbf{Robustness to semantic/map degradation:} added two studies: (A) random semantic detection drop (20\%, 40\%) and (B) section landmark removal (30\%, 50\%) with recovery analysis.
  \item \textbf{Runtime evaluation:} added repeated runtime profiling with full-pipeline throughput and component-level timings.
  \item \textbf{Experimental reproducibility details:} expanded hardware/sensor/setup details (Thorvald platform, Intel RealSense D435i, SICK Tim7 2D LiDAR, RTK-GNSS), map usage, and fixed parameters.
  \item \textbf{Figure readability:} regenerated trajectory figures with larger labels/legends, clearer line rendering, landmark/semantic-wall overlays, and start/end markers.
  \item \textbf{Text fixes:} corrected minor typographical and grammar issues raised by the reviewers.
\end{itemize}

\newpage

\section*{Associate Editor}
\paragraph{Report}
{\color{commentColor}
  (1) Experimental evaluation is difficult to interpret; fairness concern because the proposed method uses GPS while some baselines do not; clarify RPE as RPE@X meters.\\
  (2) The advantages of semantics (especially semantic walls) are not sufficiently clear methodologically and quantitatively.\\
  (3) Plot readability needs improvement.
}

\answer{to Associate Editor}{
  Thank you for the detailed guidance. We revised the manuscript along these three axes:
  \textbf{(1) Interpretability and fairness.}
  We added an \emph{AMCL+NoisyGPS} baseline and reran experiments on both traverses with three seeds.
  We now report raw APE and \emph{RPE@2m/@5m}.
  For example, in Experiment~2, SLPF achieves raw APE $=1.24\pm0.04$ m, compared to $3.38\pm0.91$ m for AMCL+NoisyGPS and $3.50\pm0.94$ m for AMCL.
  \textbf{(2) Semantic contribution clarity.}
  We clarified the semantic-wall formulation and added explicit ablations.
  In Experiment~1, removing semantic likelihood reduces row-correct fraction from $0.73$ to $0.67$ and increases cross-track error from $1.26$ m to $1.50$ m.
  \textbf{(3) Plot clarity.}
  We regenerated the affected figures with larger fonts, thicker trajectories, consistent scaling, and overlays of landmarks and semantic walls.
}

\newpage
\section*{Reviewer 1}

\paragraph{Comment 1.1}
{\color{commentColor}
  The introduction/related work underplays prior fusion methods. Clarify novelty and justify 2D LiDAR choice.
}
\answer{1.1}{
  We expanded the introduction and related-work discussion to better position prior LiDAR–vision fusion approaches.
  We clarify that our main contribution is a semantic particle-filter formulation embedding row-level structural constraints into the likelihood with adaptive GNSS weighting to address vineyard aliasing.
  We also added a discussion of the 2D versus 3D LiDAR trade-off and rationale for the selected sensing configuration.
}

\paragraph{Comment 1.2}
{\color{commentColor}
  Clarify abstract/introduction: RGB-D semantics, AMCL definition, effect sizes, LiDAR type, and navigation accuracy.
}
\answer{1.2}{
  We revised the abstract and introduction to:
  \begin{itemize}
    \item explicitly describe semantics as RGB-D driven;
    \item expand AMCL at first use;
    \item include quantitative headline improvements;
    \item state the 2D LiDAR modality;
    \item clarify operational row-level accuracy requirements motivating the chosen metrics.
  \end{itemize}
}

\paragraph{Comment 1.3}
{\color{commentColor}
  Clarify data flow, calibration, synchronisation, and map assumptions.
}
\answer{1.3}{
  We clarified the end-to-end data flow from RGB-D masks to BEV projection and semantic-LiDAR labelling.
  Fixed parameter settings are now explicitly reported.
  We state that synchronised RGB-D, LiDAR, and RTK streams are used with known camera–robot extrinsics.
  Map construction assumptions and maintenance considerations are discussed, and robustness to degraded maps is quantified.
}

\paragraph{Comment 1.4}
{\color{commentColor}
  Clarify experimental setup and baseline settings; improve qualitative figures.
}
\answer{1.4}{
  We expanded the site description (10 rows, $\sim$16.86 m average row length, $\sim$385 m$^2$ footprint) and sensor details (SICK Tim7 2D LiDAR, Intel RealSense D435i).
  For fairness:
  \begin{itemize}
    \item RTAB-Map RGB/RGBD runs use default configurations;
    \item AMCL LiDAR range truncated to 5 m for comparability with the SLPF;
    \item AMCL+NoisyGPS baseline added;
    \item multiseed evaluation used throughout.
  \end{itemize}
  Trajectory figures were regenerated with improved readability and semantic-wall overlays.
}

\paragraph{Comment 1.5}
{\color{commentColor}
  Quantify failure modes via degraded detections or map perturbations.
}
\answer{1.5}{
  We added two robustness studies on Experiment~1 (three seeds).
  \textbf{Detection Dropout:}
  Baseline raw APE is $1.00\pm0.03$ m.
  With 20\% dropout: $1.02\pm0.08$ m.
  With 40\% dropout: $1.05\pm0.09$ m.
  \textbf{Section Landmark Removal:}
  With 50\% removal, raw APE remains $1.13\pm0.13$ m and row-correct remains $0.72$.
  All runs successfully recover after exiting degraded regions.
  These results demonstrate graceful degradation and recovery behaviour.
}

\newpage
\section*{Reviewer 2}

\paragraph{Comment 2.1}
{\color{commentColor}
  Clarify semantic-wall fusion and compare to real-landmark-only formulation.
}
\answer{2.1}{
  We clarified semantic-wall construction from adjacent landmarks and detailed ray casting and class-aware likelihood formulation.
  In the ablation (Experiment~1), removing semantic likelihood reduces row-correct from $0.73$ to $0.67$ and increases cross-track from $1.26$ m to $1.50$ m.
}

\paragraph{Comment 2.2}
{\color{commentColor}
  Ensure fair comparison since the proposed method uses GPS.
}
\answer{2.2}{
  We added an AMCL+NoisyGPS fusion with a Kalman filter as an additional baseline.
  In Experiment~1, SLPF achieves $1.07\pm0.09$ m vs $1.33\pm0.46$ m for AMCL+NoisyGPS.
  In Experiment~2, SLPF achieves $1.24\pm0.04$ m vs $3.38\pm0.91$ m for AMCL+NoisyGPS.
  All methods now use identical multiseed evaluation and metrics.
}

\paragraph{Comment 2.3}
{\color{commentColor}
  Provide runtime evaluation.
}
\answer{2.3}{
  Runtime profiling (three trials, warmup excluded) yields:
  \begin{itemize}
    \item processed throughput: $13.04\pm0.21$ Hz,
    \item stride-adjusted throughput: $52.18\pm0.84$ Hz.
  \end{itemize}
  Component-level timing breakdown is included.
}

\paragraph{Comment 2.4}
{\color{commentColor}
  Minor readability and grammar issues.
}
\answer{2.4}{
  All minor issues were corrected, including figure readability, caption punctuation, and typographical errors.
}

\end{document}
